% ============================================================
% AUTOMATISCH GENERIERT aus abgabe_korpus.tex + Quellen.bib
% Skript: generate_abgabe_literatur.py
% Golden Source: abgabe_korpus.tex
% Keine manuellen Aenderungen! Bei Korrekturen:
%   1. abgabe_korpus.tex anpassen
%   2. python generate_abgabe_literatur.py ausfuehren
% ============================================================

\subsection{Kernliteratur zu Frage 1}

\textbf{Gezählte Kernliteratur}

\begin{description}
  \item[\cite{stamm2021fehlenderblick}] Der fehlende Blick auf begabte Minoritäten.
    
    In: \textit{Handbuch Begabung}.
    S.\,576--585 (10\,S.)

  \item[\cite{stamm2025vonuntennachoben}] Von unten nach oben.
    S.\,36--37, 58--62 (7\,S.)

  \item[\cite{preckel2013hochbegabung}] Hochbegabung.
    S.\,15--47 (33\,S.)

  \item[\cite{gauckreimann2021psychdiagnostik}] Psychologische Diagnostik in der Begabungs- und Begabtenförderung.
    
    In: \textit{Handbuch Begabung}.
    S.\,239--245 (7\,S.)

  \item[\cite{haag2018leistungsstanddiagnostik}] Leistungsstanddiagnostik.
    
    In: \textit{Diagnostik bei Migrantinnen und Migranten: Ein Handbuch}.
    S.\,153--165, 187--188 (15\,S.)

  \item[\cite{kellerkoller2025hellekoepfe}] Helle Köpfe mit Migrationshintergrund.
    
    In: \textit{Lichtblick für helle Köpfe}.
    S.\,76--78 (3\,S.)

  \item[\cite{baumschader2021twice}] {Twice Exceptionality} -- in zweifacher Hinsicht aussergewöhnlich.
    
    In: \textit{Handbuch Begabung}.
    S.\,588--598 (11\,S.)

  \item[\cite{koop2025herkunft}] Die Herkunft entscheidet zu oft über die Chancen.
    
    In: \textit{Was ist fair? Begabungsgerechtigkeit auf dem Prüfstand}.
    S.\,64--67 (4\,S.)

  \item[\cite{stern2025intelligenz}] Die Intelligenz kann sich im Schulalter noch verändern.
    
    \textit{Bildung Schweiz: Zeitschrift des LCH}.
    S.\,15--18 (4\,S.)

  \item[\cite{webb2020doppeldiagnosen}] Doppeldiagnosen und Fehldiagnosen bei Hochbegabung.
    S.\,87--94 (8\,S.)

  \item[\cite{warneckehauke2020bildungsgerechtigkeit}] Stärkung der Bildungsgerechtigkeit bei Underachievement, Migration und Hochbegabung.
    
    In: \textit{Begabungsförderung, Leistungsentwicklung, Bildungsgerechtigkeit -- für alle! II. Beiträge aus der Begabungsförderung}.
    S.\,242--253 (12\,S.)

  \item[\cite{kellerkoller2011erkennen}] Begabte mit Migrationshintergrund -- Erkennen und Fördern.
    S.\,1--2 (2\,S.)

  \item[\cite{muelleroppliger2021paeddiagnostik}] Pädagogische Diagnostik.
    
    In: \textit{Handbuch Begabung}.
    S.\,224--235 (12\,S.)

  \item[\cite{kappus2010migration}] Umgang mit migrationsbedingter Heterogenität.
    
    In: \textit{Alle gleich -- alle unterschiedlich}.
    S.\,63--70, 74 (9\,S.)

\end{description}

\textbf{Stützliteratur (nicht gezählt)}

\begin{description}
  \item[\cite{baudson2021wasdenken}] Was Menschen über (Hoch-)Begabung und (Hoch-)Begabte denken.
    
    In: \textit{Handbuch Begabung}.
    S.\,115--128.

  \item[\cite{muelleroppliger2021begabungsmodelle}] Begabungsmodelle.
    
    In: \textit{Handbuch Begabung}.
    S.\,204--219.

  \item[\cite{baudson2025besserfinden}] Wie wir Begabte besser finden können.
    
    In: \textit{Was ist fair? Begabungsgerechtigkeit auf dem Prüfstand}.
    S.\,35--40.

  \item[\cite{trautmann2016einfuehrung}] Einführung in die Hochbegabtenpädagogik.
    S.\,34--41.

  \item[\cite{lehwald2017motivation}] Motivation trifft Begabung.
    S.\,78--92.

  \item[\cite{grossenbacher2014integrative}] Talent und Begabung in der Volksschule der deutschsprachigen Schweiz.
    
    In: \textit{Handbuch Talententwicklung}.
    S.\,317--325.

  \item[\cite{erzinger2023pisa}] PISA 2022.
    S.\,45--48.

  \item[\cite{bfs2022migration}] Migration und Integration.
    S.\,34--35.

  \item[\cite{dvs2025bbf}] Integrative Begabungs- und Begabtenförderung.
    Online.

  \item[\cite{maehler2018diagnostik}] Diagnostik bei Migrantinnen und Migranten.
    S.\,153--165,187f..

\end{description}

\textit{Umfang Frage 1: 137\,Seiten gezählte Kernliteratur; Stützliteratur nicht eingerechnet.}


\subsection{Kernliteratur zu Frage 2}

\textbf{Gezählte Kernliteratur}

\begin{description}
  \item[\cite{saegesserwyss2022grafinkrahmenmodell}] Grafomotorik und schulische Inklusion.
    
    In: \textit{Gesellschaft und Diversit\"at}.
    S.\,2--3, 6--11 (7\,S.)

  \item[\cite{nottbusch2017graphomotorik}] Graphomotorik.
    
    In: \textit{Forschungshandbuch empirische Schreibdidaktik}.
    S.\,128--136 (9\,S.)

  \item[\cite{hurschler2020handschriftbeurteilung}] Differenzierende Beurteilung der Handschrift.
    
    \textit{leseforum.ch: Online-Plattform für Literalität}.
    S.\,1--24 (24\,S.)

  \item[\cite{gold2018lesenkannmanlernen}] Lesen kann man lernen.
    S.\,50--53, 62, 67--88 (25\,S.)

  \item[\cite{lehwald2017motivation}] Motivation trifft Begabung.
    S.\,70--72, 84--89 (9\,S.)

  \item[\cite{greiten2021underachievement}] Underachievement.
    
    In: \textit{Handbuch Begabung}.
    S.\,546--553 (8\,S.)

\end{description}

\textbf{Stützliteratur (nicht gezählt)}

\begin{description}
  \item[\cite{baumschader2021twice}] {Twice Exceptionality} -- in zweifacher Hinsicht aussergewöhnlich.
    
    In: \textit{Handbuch Begabung}.
    S.\,588--598.

  \item[\cite{warneckehauke2020bildungsgerechtigkeit}] Stärkung der Bildungsgerechtigkeit bei Underachievement, Migration und Hochbegabung.
    
    In: \textit{Begabungsförderung, Leistungsentwicklung, Bildungsgerechtigkeit -- für alle! II. Beiträge aus der Begabungsförderung}.
    S.\,247--252.

  \item[\cite{hoyer2013begabung}] Begabung.
    S.\,111--113.

  \item[\cite{sturm2016graphomotorik}] Zur Bedeutung der graphomotorischen Prozesse beim Schreiben(lernen).
    
    In: \textit{Facetten grundschulpädagogischer und -didaktischer Forschung}.
    S.\,183--198.

\end{description}

\textit{Umfang Frage 2: 82\,Seiten gezählte Kernliteratur; Stützliteratur nicht eingerechnet.}


\subsection{Kernliteratur zu Frage 3}

\textbf{Gezählte Kernliteratur}

\begin{description}
  \item[\cite{grossrieder2010anerkennung}] Gleiches und Unterschiedliches anerkennen.
    
    In: \textit{Alle gleich -- alle unterschiedlich}.
    S.\,88--94 (7\,S.)

  \item[\cite{baudson2021wasdenken}] Was Menschen über (Hoch-)Begabung und (Hoch-)Begabte denken.
    
    In: \textit{Handbuch Begabung}.
    S.\,115--129 (15\,S.)

  \item[\cite{kuhl2021begabungbildungbeziehung}] Begabung, Bildung und Beziehung aus persönlichkeitspsychologischer Sicht.
    
    In: \textit{Handbuch Begabung}.
    S.\,185--201 (17\,S.)

  \item[\cite{wagener2021bfhemmendfoerdernd}] Begabungsfördernde und hemmende Wirkungsmechanismen im Klassenkontext.
    
    In: \textit{Handbuch Begabung}.
    S.\,418--424 (7\,S.)

  \item[\cite{behrensen2019inklusive}] Inklusive Begabungsförderung -- eine Antwort auf Bildungsherausforderungen in Zeiten von Fluchtmigration.
    
    In: \textit{Begabungsförderung und Professionalisierung}.
    S.\,86--98 (13\,S.)

  \item[\cite{kuhl2019diversitaet}] Diversität und Persönlichkeit.
    
    In: \textit{Begabungsförderung und Professionalisierung}.
    S.\,35--54 (20\,S.)

  \item[\cite{boosnuenning2022interethnisch}] Interethnische Freundschaften unter den Bedingungen sozialräumlicher und institutioneller Segregation.
    S.\,51--65 (15\,S.)

  \item[\cite{kesselshannover2015gleichaltrige}] Gleichaltrige.
    
    In: \textit{Pädagogische Psychologie}.
    S.\,288 (1\,S.)

  \item[\cite{tschoppgruetterbuholzer2022intergruppenkontakt}] Soziale Teilhabe durch Intergruppenkontakt.
    S.\,35--40 (6\,S.)

  \item[\cite{baumert2022freundschaftwerte}] {Man kann nicht leben, wenn man keine Freunde hat}.
    S.\,40--49 (10\,S.)

  \item[\cite{evers2025stress}] Stress als Dämpfer für die kognitive Entwicklung und die Entfaltung von Begabung.
    
    In: \textit{Was ist fair? Begabungsgerechtigkeit auf dem Prüfstand}.
    S.\,21--27 (8\,S.)

  \item[\cite{lehwald2017motivation}] Motivation trifft Begabung.
    S.\,57--63 (7\,S.)

\end{description}

\textbf{Stützliteratur (nicht gezählt)}

\begin{description}
  \item[\cite{baumschader2021twice}] {Twice Exceptionality} -- in zweifacher Hinsicht aussergewöhnlich.
    
    In: \textit{Handbuch Begabung}.
    S.\,588--598.

  \item[\cite{kappus2010migration}] Umgang mit migrationsbedingter Heterogenität.
    
    In: \textit{Alle gleich -- alle unterschiedlich}.
    S.\,63--70, 74.

  \item[\cite{gysinscherzinger2022freundschaftenwertvoll}] Freundschaften machen das Leben wertvoll.
    S.\,8--12.

\end{description}

\textit{Umfang Frage 3: 126\,Seiten gezählte Kernliteratur; Stützliteratur nicht eingerechnet.}


\subsection{Kernliteratur zu Frage 4}

\textbf{Gezählte Kernliteratur}

\begin{description}
  \item[\cite{buholzerkummerwyss2010reaktionen}] Zur Einführung.
    
    In: \textit{Alle gleich -- alle unterschiedlich}.
    S.\,78--85 (8\,S.)

  \item[\cite{muelleroppliger2021plurale}] Plurale Gesellschaft, Inklusion und Bildungsgerechtigkeit.
    
    In: \textit{Handbuch Begabung}.
    S.\,32--42 (11\,S.)

  \item[\cite{muelleroppliger2021begabungsmodelle}] Begabungsmodelle.
    
    In: \textit{Handbuch Begabung}.
    S.\,204--219 (16\,S.)

  \item[\cite{reisrenzullimueller2021sem}] Das Schoolwide Enrichment Model (SEM).
    
    In: \textit{Handbuch Begabung}.
    S.\,333--345 (13\,S.)

  \item[\cite{muelleroppliger2021paeddiagnostik}] Pädagogische Diagnostik.
    
    In: \textit{Handbuch Begabung}.
    S.\,224--235 (0\,S.)

  \item[\cite{weigand2021separativ}] Separativ oder integrativ? Inklusive Begabungs- und Begabtenförderung.
    
    In: \textit{Handbuch Begabung}.
    S.\,290--298 (9\,S.)

  \item[\cite{muelleroppliger2021adaptive}] Adaptive Lernarchitekturen.
    
    In: \textit{Handbuch Begabung}.
    S.\,374--385 (12\,S.)

  \item[\cite{schulteterhardt2020potenzialentwicklung}] Individuelle Potenzialentwicklung durch stärkenorientierte Lernarchitekturen.
    
    In: \textit{Begabungsförderung, Leistungsentwicklung, Bildungsgerechtigkeit -- für alle! Beiträge aus der Begabungsförderung}.
    S.\,264--267 (4\,S.)

  \item[\cite{lehwald2017motivation}] Motivation trifft Begabung.
    S.\,151--158 (8\,S.)

  \item[\cite{hoyer2013begabung}] Begabung.
    S.\,94--99 (6\,S.)

  \item[\cite{sedmak2021bildungsgerechtigkeit}] Begabtenförderung und Bildungsgerechtigkeit.
    
    In: \textit{Handbuch Begabung}.
    S.\,65--75 (11\,S.)

  \item[\cite{grossenbacher2014integrative}] Talent und Begabung in der Volksschule der deutschsprachigen Schweiz.
    
    In: \textit{Handbuch Talententwicklung}.
    S.\,317--325 (9\,S.)

\end{description}

\textbf{Stützliteratur (nicht gezählt)}

\begin{description}
  \item[\cite{horvath2021elite}] Elite, Begabung und soziale Ungleichheit.
    
    In: \textit{Handbuch Begabung}.
    S.\,77--87.

  \item[\cite{weigand2021person}] Begabung, Bildung und Person.
    
    In: \textit{Handbuch Begabung}.
    S.\,46--59.

  \item[\cite{baudson2021wasdenken}] Was Menschen über (Hoch-)Begabung und (Hoch-)Begabte denken.
    
    In: \textit{Handbuch Begabung}.
    S.\,119.

  \item[\cite{dvs2025bbf}] Integrative Begabungs- und Begabtenförderung.
    Online.

\end{description}

\textit{Umfang Frage 4: 107\,Seiten gezählte Kernliteratur; Stützliteratur nicht eingerechnet.}


\subsection{Kernliteratur zu Frage 5}

\textbf{Gezählte Kernliteratur}

\begin{description}
  \item[\cite{weigand2021person}] Begabung, Bildung und Person.
    
    In: \textit{Handbuch Begabung}.
    S.\,46--59 (14\,S.)

  \item[\cite{horvath2021elite}] Elite, Begabung und soziale Ungleichheit.
    
    In: \textit{Handbuch Begabung}.
    S.\,77--87 (11\,S.)

  \item[\cite{muellerboeschschaffnermenn2021udl}] Inklusiver Unterricht.
    
    In: \textit{Inklusive P\"adagogik und Didaktik}.
    S.\,94--116 (23\,S.)

  \item[\cite{macha2019gender}] Gender- und diversitätssensible Begabungsförderung und Professionalisierung.
    
    In: \textit{Begabungsförderung und Professionalisierung}.
    S.\,160--171 (12\,S.)

  \item[\cite{groschefussangelgraesel2020kokonstruktion}] Kokonstruktive Kooperation zwischen Lehrkräften.
    
    \textit{Zeitschrift f\"ur P\"adagogik}.
    S.\,463--467, 469--472 (9\,S.)

  \item[\cite{nguyensliwka2021massnahmen}] Schulische Massnahmen und Kompetenzen von Lehrpersonen für die Begabungsförderung.
    
    In: \textit{Handbuch Begabung}.
    S.\,348--356 (9\,S.)

  \item[\cite{widmerwolf2018multiprofessionell}] Kooperation in multiprofessionellen Teams an inklusiven Schulen.
    
    In: \textit{Handbuch schulische Inklusion}.
    S.\,299--303, 308--309 (7\,S.)

  \item[\cite{kosoroklabhart2021voneltern}] Von Eltern mit Migrationshintergrund lernen.
    S.\,14--15, 38--39 (4\,S.)

  \item[\cite{baudson2025besserfinden}] Wie wir Begabte besser finden können.
    
    In: \textit{Was ist fair? Begabungsgerechtigkeit auf dem Prüfstand}.
    S.\,35--40 (6\,S.)

\end{description}

\textbf{Stützliteratur (nicht gezählt)}

\begin{description}
  \item[\cite{muelleroppliger2021paeddiagnostik}] Pädagogische Diagnostik.
    
    In: \textit{Handbuch Begabung}.
    S.\,224--238.

  \item[\cite{reisrenzullimueller2021sem}] Das Schoolwide Enrichment Model (SEM).
    
    In: \textit{Handbuch Begabung}.
    S.\,333--347.

  \item[\cite{weigand2021separativ}] Separativ oder integrativ? Inklusive Begabungs- und Begabtenförderung.
    
    In: \textit{Handbuch Begabung}.
    S.\,290--298.

  \item[\cite{muelleroppliger2021plurale}] Plurale Gesellschaft, Inklusion und Bildungsgerechtigkeit.
    
    In: \textit{Handbuch Begabung}.
    S.\,32--45.

  \item[\cite{sedmak2021bildungsgerechtigkeit}] Begabtenförderung und Bildungsgerechtigkeit.
    
    In: \textit{Handbuch Begabung}.
    S.\,65--75.

  \item[\cite{muelleroppliger2021adaptive}] Adaptive Lernarchitekturen.
    
    In: \textit{Handbuch Begabung}.
    S.\,384--385.

  \item[\cite{wagener2021bfhemmendfoerdernd}] Begabungsfördernde und hemmende Wirkungsmechanismen im Klassenkontext.
    
    In: \textit{Handbuch Begabung}.
    S.\,418--424.

  \item[\cite{kummerwyss2017kooperativunterrichten}] Kooperativ unterrichten.
    
    In: \textit{Alle gleich -- alle unterschiedlich! Zum Umgang mit Heterogenität in Schule und Unterricht}.
    S.\,151--154, 156--160.

\end{description}

\textit{Umfang Frage 5: 95\,Seiten gezählte Kernliteratur; Stützliteratur nicht eingerechnet.}

